\section{Motivation and research questions}\label{sec:section-01}

A trading research system must answer two questions that are easy to confuse: whether its calculations represent the stated strategy, and whether that strategy has useful economic performance. Correct accounting is necessary for the second question, but it cannot answer it. Likewise, a profitable curve cannot validate data that were unavailable at the simulated decision time. This study follows the entire chain from public observations to predictions, constrained targets, simulated fills, persisted account state and research claims. The output is an executable local paper-trading product and a versioned empirical record. No external broker execution is part of the experiment.

The starting repository contained classifiers, a strategy zoo, an ensemble, a dashboard and saved backtests. Earlier audits identified accounting and timing defects. Subsequent repairs produced a sealed adaptive experiment and a short technical report. We preserve that experiment and examine its full prediction and ledger records instead of inferring behavior from headline accuracy or isolated screenshots. This follow-up adds class and instrument diagnostics, controlled refits, transfer tests, execution sensitivities and operational evidence. Previously inspected historical observations are explicitly exploratory. A new analysis of an old holdout does not restore its independence.

Our first question is whether causal forecasts improve probability quality and directional classification over a trailing class-frequency estimate and simple forecasting controls. The primary predictive endpoint is the paired, date-averaged log-loss difference against the prior. Accuracy, balanced accuracy, class F1 and multiclass Brier score explain different aspects of the result. A correct class prediction and a profitable trade are separate events, so no accuracy threshold is used as a proof of alpha.

The second question is whether a hybrid adds net economic value over its components and comparable simple policies when every strategy must finance its holdings through one account. The primary economic endpoint is the daily net-return difference from a costed, risk-constrained equal-weight account on matching dates. We distinguish this account from a descriptive frictionless SPY reference. The third question concerns where errors occur: classes, instruments, calendar blocks, observed volatility conditions, disagreement between experts and confident mistakes. Such conditional descriptions identify hypotheses; they do not by themselves establish a causal mechanism.

The fourth question is sensitivity to features, calibration, gate construction, training history, refit cadence and execution assumptions. Controlled comparisons change a specified factor while preserving the eligible information set and outcome definition. The fifth question concerns operation: whether the same policy can run from fresh inputs with persistent capital, restart after faults, refuse invalid state and produce durable evidence on its registered schedule. These questions define separate engineering, historical-evidence, prospective-operation and manuscript decisions. Broker readiness is a further decision outside the authorized scope.
